% Define the document type and font size
\documentclass[12pt]{article}

\usepackage[spanish]{babel} % Package for the Spanish language
\usepackage[utf8]{inputenc} % Package for UTF-8 character encoding
\usepackage{geometry} % Package to configure document margins
\usepackage{graphicx} % Package to include graphics
\usepackage{listings} % Package to include source code in the document
\usepackage{xcolor} % Package to define colors
\usepackage{fancyhdr} % Package to customize headers and footers
\usepackage{amsmath} % Package for advanced mathematical symbols and environments
\usepackage{amssymb} % Package for additional mathematical symbols
\usepackage[hidelinks]{hyperref} % Package for hyperlinks in index and references (Keep this as the last package loaded)

\setlength{\parskip}{1em} % Space between paragraphs
\setlength{\parindent}{0.0in} % Indentation for paragraphs

% Definition of colors for source code
\definecolor{COLOR_COMMENTS}{HTML}{999999}
\definecolor{COLOR_KEYWORDS}{HTML}{37a5a6}
\definecolor{COLOR_IDENTIFIERS}{HTML}{ea76cb}
\definecolor{COLOR_STRINGS}{HTML}{ff7858}

% Configuration of the listings environment to display source code
\lstset{
frame=shadowbox,
aboveskip=3mm,
belowskip=3mm,
xleftmargin=10mm,
xrightmargin=10mm,
showstringspaces=false,
columns=flexible,
basicstyle={\footnotesize\ttfamily},
numbers=left,
numberstyle=\tiny\color{COLOR_COMMENTS},
keywordstyle=\color{COLOR_KEYWORDS},
stringstyle=\color{COLOR_STRINGS},
commentstyle=\color{COLOR_COMMENTS},
breaklines=true,
breakatwhitespace=true,
tabsize=4,
}

% Configuration of document margins
\geometry{
    a4paper,
    left=25mm,
    right=25mm,
    top=25mm,
    bottom=25mm
}
\graphicspath{{C:/Users/mishi/projects/LaTeX}} % Path for images

\title{Aproximación de Funciones mediante Redes de Base Radial (RBF)} % Document title
\author{María Laura Reina y Victoria Ruza} % Document author
\date{\today} % Document date

\begin{document}

    \begin{titlepage}
        \centering
        {\includegraphics[width=0.2\textwidth]{logoBlack}\par}
        {\large {Universidad Central de Venezuela}\par}
        {\large {Facultad de Ciencias}\par}
        {\large {Escuela de Computación}\par}
        {\large {Cálculo Científico}\par}
        \vfill
        {\Huge {Aproximación de Funciones mediante Redes de Base Radial (RBF)}\par}
        \vspace{0.5cm}
        {\Large {Taller 6}\par}
        \vfill
        {\large María Laura Reina y Victoria Ruza\par}
        \vspace{0.5cm}
        {\large \today\par}
    \end{titlepage}

    \pagenumbering{gobble}
    \tableofcontents
    \newpage
    \pagenumbering{arabic}

    \section{Referencias}
    \begin{itemize}
        \item Radial Basis Interpolation by jessebett. (n.d.). \url{https://www.cs.toronto.edu/~jessebett/projects/radial-basis-functions/}
        \item Mongillo, M., 2011. Choosing Basis Functions and Shape Parameters for Radial Basis Function Methods, Available at: \url{http://www.siam.org/students/siuro/vol4/S01084.pdf}
    \end{itemize}

\end{document}